\documentclass[review]{elsarticle}
\usepackage{lineno}
\usepackage{amssymb}
\usepackage{amsmath}
\usepackage{booktabs}
\usepackage{siunitx}
\usepackage{enumitem}
\usepackage{graphicx}
\usepackage{hyperref}

\journal{Engineering Applications of Artificial Intelligence}
\date{}

\begin{document}
\linenumbers
\begin{frontmatter}
\title{Impact of Quantization and Structured Pruning on Battery SOH Prediction Models across Architectures}

\author{Anonymous}
\affiliation{organization={Anonymous}}

\begin{abstract}
\begin{itemize}[leftmargin=*]
\item Core objective: investigate how pruning, quantization, and knowledge distillation affect SOH prediction quality and embedded deployment cost across LSTM, GRU, CNN, and TCN architectures.
\item Study principle: architectures are trained under a harmonized protocol first, then optimized under controlled compression settings to enable fair cross-architecture comparison.
\item Evaluation focus: MAE/RMSE on SOH prediction and embedded KPIs (flash, RAM, latency, energy).
\item Expected outcome: architecture-specific design guidance for resource-constrained BMS implementations.
\end{itemize}
\end{abstract}

\begin{keyword}
Battery SOH prediction \sep Structured pruning \sep Quantization \sep Knowledge distillation \sep Embedded AI \sep LSTM \sep GRU \sep CNN \sep TCN \sep BMS
\end{keyword}
\end{frontmatter}

\section{Introduction}
\begin{itemize}[leftmargin=*]
\item Battery management systems require accurate SOH prediction under strict constraints on flash, RAM, latency, and energy.
\item Many published studies emphasize predictive accuracy but report only limited evidence for embedded deployment trade-offs.
\item This paper focuses on architecture-dependent optimization behavior and asks whether different model families benefit equally from pruning, quantization, and distillation.
\item The target is a comparable SOH-accuracy range across architectures before compression, followed by systematic compression analysis.
\item Figure 1 (planned, position in this section): high-level study concept showing the flow from baseline models to compression methods and final embedded evaluation.
\end{itemize}

\section{Related Work and Research Gap}
\begin{itemize}[leftmargin=*]
\item Existing SOH literature covers recurrent and convolutional sequence models, but direct cross-architecture compression comparisons are still limited.
\item Most compression studies are model-centric and do not combine battery-domain evaluation with embedded KPIs in one unified protocol.
\item The open gap is a controlled cross-architecture analysis that links compression decisions directly to BMS-relevant engineering costs.
\item Figure 2 (planned, position in this section): literature positioning map contrasting accuracy-centric studies, compression-centric studies, and the proposed integrated approach.
\end{itemize}

\section{Methodology}
\subsection{Comparison Objective and Fairness Principle}
\begin{itemize}[leftmargin=*]
\item Step 1: establish strong baseline models per architecture in a comparable SOH-accuracy band.
\item Step 2: apply optimization methods under controlled settings and evaluate accuracy-cost shifts.
\item Step 3: compare architecture-specific robustness and derive deployment-oriented recommendations.
\item Fairness rule: model sizes may differ when required to reach comparable baseline quality, but data split, metric definition, and benchmark protocol remain fixed.
\end{itemize}

\subsection{Data, Target, and Split Strategy}
\begin{itemize}[leftmargin=*]
\item Prediction target: battery SOH over aging trajectories.
\item Feature basis from current code: Voltage, Current, Temperature, EFC, and Q\_c.
\item Feature preprocessing: hourly aggregation into mean, std, min, and max channels.
\item Common partitioning: train cells C01/C03/C05/C09/C13/C17/C25/C27/C29; validation cells C07/C15/C21; test cells C11/C19/C23.
\item Figure 3 (planned, position in this subsection): SOH trajectories of all cells with color-coded train/validation/test assignment.
\end{itemize}
\noindent This split figure documents partition coverage and helps verify that all sets contain representative SOH evolution ranges.

\subsection{Architecture Portfolio}
\begin{itemize}[leftmargin=*]
\item LSTM reference model: version 0.1.2.3, hidden size 192, 3 recurrent layers.
\item GRU reference model: version 0.3.1.1, hidden size 256, 4 recurrent layers.
\item TCN reference model: version 0.2.2.1, hidden size 96 with causal dilation stack.
\item CNN reference model: version 0.4.1.1, hidden size 224 with dilated causal blocks.
\item Figure 4 (planned, position in this subsection): architecture schematic panel with one compact diagram per model family.
\end{itemize}

\begin{table*}[t]
\centering
\caption{Architecture comparison for the SOH optimization study (baseline FP32 models).}
\label{tab:arch_portfolio}
\small
\begin{tabular}{llcccc}
\toprule
Family & Version & Hidden Size & Main Depth & Parameters & Key Structure \\
\midrule
LSTM & 0.1.2.3 & 192 & 3 LSTM layers & 1,041,185 & LSTM-Hybrid-Seq2Seq with residual MLP head \\
GRU  & 0.3.1.1 & 256 & 4 GRU layers  & 1,919,617 & GRU-Hybrid-Seq2Seq with residual MLP head \\
TCN  & 0.2.2.1 & 96  & 4 blocks, dilations 1/2/4/8/16 & 266,977 & Causal TCN Seq2Seq \\
CNN  & 0.4.1.1 & 224 & 4 dilated causal blocks & 2,056,929 & Dilated causal CNN Seq2Seq \\
\bottomrule
\end{tabular}
\end{table*}

\subsection{Training and Evaluation Protocol}
\begin{itemize}[leftmargin=*]
\item Training setup is harmonized across architectures (seed 42, seq2seq objective, smoothness regularization, cosine warm restarts, early stopping, gradient clipping).
\item Main predictive metrics are MAE and RMSE; optional analyses can include EOL-oriented indicators.
\item Current practical workflow in the project: one main run per selected model version, with model acceptance based on test MAE in a practically usable range.
\item Figure 5 (planned, position in this subsection): training/evaluation flow chart from data loading to baseline model selection.
\end{itemize}

\subsection{Model Optimization Methods}
\begin{itemize}[leftmargin=*]
\item Pruning: structured pruning is applied in percentage steps to quantify architecture-specific sensitivity and breakpoints.
\item Quantization: precision is reduced from FP32 to lower-bit formats to evaluate memory and runtime savings versus accuracy loss.
\item Knowledge distillation: teacher-student training is integrated to recover predictive quality of compressed models.
\item Combined settings: selected model variants are evaluated with joint optimization (for example pruning + quantization, with optional distillation support).
\item Figure 6 (planned, position in this subsection): optimization pipeline diagram showing baseline model, individual optimization branches, and combined branch.
\end{itemize}

\subsection{Embedded Benchmark Protocol}
\begin{itemize}[leftmargin=*]
\item All variants are benchmarked with the same target platform and firmware toolchain.
\item Inference is performed with batch size 1 and fixed warm-up, because deployment is stateful and sample-by-sample in online BMS operation.
\item Embedded KPIs: MAE, flash, RAM, latency (mean and std), and energy.
\item Quantization path in current code: dynamic INT8 for LSTM/GRU, static FX INT8 for CNN, and stabilized head-focused static FX for TCN.
\item Figure 7 (planned, position in this subsection): benchmark setup including host pipeline, firmware execution, and metric logging flow.
\end{itemize}

\section{Results}
\subsection{Baseline Model Comparison}
\begin{itemize}[leftmargin=*]
\item Report baseline MAE/RMSE per architecture on the shared test set.
\item Compare baseline model sizes and baseline embedded costs before any compression.
\item Figure 8 (planned, position in this subsection): grouped bars for baseline MAE/RMSE and baseline footprint.
\end{itemize}

\subsection{Pruning Results}
\begin{itemize}[leftmargin=*]
\item Report MAE/RMSE trends over pruning levels per architecture.
\item Report flash/RAM/latency/energy changes over pruning levels.
\item Identify pruning thresholds where performance leaves the acceptable operating range.
\item Figure 9 (planned, position in this subsection): pruning sensitivity curves and per-architecture Pareto view.
\end{itemize}

\subsection{Quantization Results}
\begin{itemize}[leftmargin=*]
\item Compare FP32 and lower precision variants per architecture.
\item Quantify trade-offs between prediction quality and embedded resource gains.
\item Distinguish robust quantization cases from degradation-prone cases.
\item Figure 10 (planned, position in this subsection): precision-vs-MAE and precision-vs-resource multi-panel plot.
\end{itemize}

\subsection{Distillation and Combined Optimization Results}
\begin{itemize}[leftmargin=*]
\item Evaluate whether distillation recovers accuracy for compressed student models.
\item Compare single-method and combined-method variants within one decision space.
\item Analyze additive versus interaction effects between optimization methods.
\item Figure 11 (planned, position in this subsection): combined optimization scatter with non-dominated configurations.
\end{itemize}

\subsection{Cross-Architecture Decision View}
\begin{itemize}[leftmargin=*]
\item Summarize all accepted variants in architecture-wise and global Pareto fronts.
\item Highlight architecture-dependent operating regions for low-memory, low-latency, and low-energy constraints.
\item Figure 12 (planned, position in this subsection): final decision map for architecture and optimization selection.
\end{itemize}

\section{Discussion}
\begin{itemize}[leftmargin=*]
\item Interpret why architectures react differently to identical optimization settings.
\item Discuss where higher model complexity is justified by robustness gains and where lightweight structures dominate.
\item Relate results to practical BMS deployment constraints and integration risks.
\item Clarify methodological limits (data regime, target platform scope, and configuration breadth).
\item Figure 13 (planned, position in this section): qualitative decision matrix linking constraints to recommended architecture-method combinations.
\end{itemize}

\section{Conclusion and Outlook}
\begin{itemize}[leftmargin=*]
\item The study provides a controlled architecture-level comparison of pruning, quantization, and distillation for SOH prediction.
\item The central contribution is a reproducible accuracy-cost analysis that directly supports embedded design decisions.
\item Future work extends the study to additional chemistries, broader operating regimes, and tighter hardware-in-the-loop integration.
\end{itemize}

\section*{Figure Directory (Working Draft)}
\begin{itemize}[leftmargin=*]
\item Figure 1: Study concept from baseline training to compression and embedded evaluation.
\item Figure 2: Literature positioning map and research gap framing.
\item Figure 3: SOH trajectories with train/validation/test split highlighting.
\item Figure 4: Architecture panel for LSTM, GRU, CNN, and TCN.
\item Figure 5: Training and baseline-selection workflow.
\item Figure 6: Optimization workflow (pruning, quantization, distillation, combined branch).
\item Figure 7: Embedded benchmark setup and metric acquisition flow.
\item Figure 8: Baseline MAE/RMSE and footprint comparison.
\item Figure 9: Pruning sensitivity and pruning Pareto view.
\item Figure 10: Quantization sensitivity and precision-resource trade-off.
\item Figure 11: Distillation and combined optimization scatter with non-dominated points.
\item Figure 12: Cross-architecture decision map.
\item Figure 13: Discussion matrix linking deployment constraints to model-method choices.
\end{itemize}

\section*{Data availability (placeholder)}
\begin{itemize}[leftmargin=*]
\item Dataset DOI, code repository links, and benchmark scripts will be inserted in the final manuscript version.
\end{itemize}

\bibliographystyle{elsarticle-num}
\bibliography{bib/paper2_socsoh_fr}
\end{document}
